                                IMC 2013, Blagoevgrad, Bulgaria
                                             Day 2, August 9, 2013



Problem 1. Let z be a complex number with |z + 1| > 2. Prove that |z 3 + 1| > 1.
                                (Proposed by Walther Janous and Gerhard Kirchner, Innsbruck)

Solution. Since z 3 + 1 = (z + 1)(z 2 − z + 1), it suffices to prove that |z 2 − z + 1| ≥ 12 .
   Assume that z + 1 = reϕi , where r = |z + 1| > 2, and ϕ = arg(z + 1) is some real number. Then

                           z 2 − z + 1 = (reϕi − 1)2 − (reϕi − 1) + 1 = r 2 e2ϕi − 3reϕi + 3,

and

                        |z 2 − z + 1|2 = r 2 e2ϕi − 3reϕi + 3) r 2 e−2ϕi − 3re−ϕi + 3) =
                                       = r 4 + 9r 2 + 9 − (6r 3 + 18r) cos ϕ + 6r 2 cos 2ϕ =
                                       = r 4 + 9r 2 + 9 − (6r 3 + 18r) cos ϕ + 6r 2 (2 cos2 ϕ − 1) =
                                                                2
                                                         r2 + 3
                                             
                                                                      1                    1   1
                                       = 12 r cos ϕ −               + (r 2 − 3)2 > 0 + = .
                                                            4         4                    4   4
This finishes the proof.


Problem 2. Let p and q be relatively prime positive integers. Prove that
                            pq−1
                                      k  k  (
                             X         p
                                          +
                                            q      0 if pq is even,
                                 (−1)          =                                                       (∗)
                             k=0
                                                   1 if pq is odd.

(Here ⌊x⌋ denotes the integer part of x.)
                                       (Proposed by Alexander Bolbot, State University, Novosibirsk)

Solution.Suppose
              k  first that pq is even (which implies that p and q have opposite parities), and let
          k
                +
           p        q
ak = (−1)          . We show that ak + apq−1−k = 0, so the terms on the left-and side of (∗) cancel out
in pairs.
    For every positive integer k we have kp + p−1−k     = p−1
                                              
                                                  p         p
                                                              , hence
                                               
    k    pq − 1 − k     k    k     pq − 1 − k     −1 − k      pq − 1 p − 1
      +              =    −     +             −             =       −      = q−1
    p        p          p    p         p            p            p     p
and similarly                                      
                                       pq − 1 − k       k
                                                    +     = p − 1.
                                            q           q
Since p and q have opposite parities, it follows that kp + kq and pq−1−k
                                                                      pq−1−k 
                                                                     p
                                                                          +     q
                                                                                      have opposite
parities and therefore apq−1−k = −ak .
   Now suppose that pq is odd. For every index k, denote by pk and qk the remainders of k modulo p
and q, respectively. (I.e., 0 ≤ pk < p and 0 ≤ qk < q such that k ≡ pk (mod p) and k ≡ qk (mod q).)
   Notice that
                                    
                 k        k        k       k
                     +        ≡p      +q      = (k − pk ) + (k − qk ) ≡ pk + qk (mod 2).
                 p        q        p       q

                                                             1
   Since p and q are co-prime, by the Chinese remainder theorem the map k 7→ (pk , qk ) is a bijection
between the sets {0, 1, . . . , pq − 1} and {0, 1, . . . , p − 1} × {0, 1, . . . , q − 1}. Hence
       pq−1      k   k  pq−1                  p−1 q−1                   p−1
                                                                                        !    q−1
                                                                                                      !
                     +
        X                        X                X      X                 X                 X
           (−1) p       q =          (−1)pk +qk =            (−1)i+j =            (−1)i ·        (−1)j = 1.
           k=0                        k=0                  i=0 j=0                             i=0               j=0




Problem 3. Suppose that v1 , . . . , vd are unit vectors in Rd . Prove that there exists a unit vector u
such that                                                  √
                                             |u · vi | ≤ 1/ d
for i = 1, 2, . . . , d.
    (Here · denotes the usual scalar product on Rd .)
                                                 (Proposed by Tomasz Tkocz, University of Warwick)

Solution. If v1 , . . . , vd are linearly dependent then we can simply take a unit vector u perpendicular to
span(v1 , . . . , vd ). So assume that v1 , . . . , vd are linearly independent. Let w1 , . . . , wd be the dual basis
of (v1 , . . . , vd ), i.e.                           (
                                                        1 if i = j
                                   wi · vj = δij =                     for 1 ≤ i, j ≤ d.
                                                        0 if i 6= j
From wi · vi = 1 we have |wi | ≥ 1.
   For every sequence ε = (ε1 , . . . , εd ) ∈ {+1, −1}d of signs define uε = di=1 εi wi . Then we have
                                                                             P

                                      d
                                      X                          d
                                                                 X
                      |uε · vk | =          εi (wi · vk ) =              εi δik = |εk | = 1 for k = 1, . . . , d.
                                      i=1                        i=1

   Now estimate the average of |uε |2 .
                                                   d
                                                                 !           d
                                                                                           !
     1       X               1          X          X                         X
                    |uε |2 = d                           εi wi       ·             εj wj       =
     2d           n
                            2
          ε∈{+1,−1}                  ε∈{+1,−1}n    i=1                       j=1
                                                                                     
                              d X
                                d                                                              d X
                                                                                                 d                         d
                              X                   1               X                            X                           X
                           =         (wi · wj )  d                            εi εj  =                 (wi · wj )δij =         |wi |2 ≥ d.
                                                 2
                             i=1 j=1                        ε∈{+1,−1}n                         i=1 j=1                     i=1

                                                                                                             uε
It follows that there is a ε such that |uε |2 ≥ d. For that ε the vector u =                                      satisfies the conditions.
                                                                                                            |uε |


Problem 4. Does there exist an infinite set M consisting of positive integers such that for any
a, b ∈ M, with a < b, the sum a + b is square-free?
    (A positive integer is called square-free if no perfect square greater than 1 divides it.)
                                            (Proposed by Fedor Petrov, St. Petersburg State University)

Solution. The answer is yes. We construct an infinite sequence 1 = n1 < 2 = n2 < n3 < . . . so that
ni + nj is square-free for all i < j. Suppose that we already have some numbers n1 < . . . < nk (k ≥ 2),
which satisfy this condition and find a suitable number nk+1 to be the next element of the sequence.
    We will choose nk+1 of the form nk+1 = 1 + Mx, with M = ((n1 + . . . + nk + 2k)!)2 and some positive
integer x. For i = 1, 2, . . . , k we have ni + nk+1 = 1 + Mx + ni = (1 + ni )mi , where mi and M are
co-prime, so any perfect square dividing 1 + Mx + ni is co-prime with M.


                                                                         2
     In order to find a suitable x, take a large N and consider the values x = 1, 2, . . . , N. If a value
1 ≤ x ≤ N is not suitable, this means that there is an index 1 ≤ i ≤ k and some prime p such that
p2 |1 + Mx + ni . For p ≤ p2k this is impossible because p|M. Moreover, we also have p2 ≤ 1 + Mx + ni <
M(N + 1), so 2k < p < M(N + 1).
     For any fixed i and p, the values for x for which p2 |1 + Mx + ni form an arithmetic progression
                                                     N
with difference p2 . Therefore, there are at most 2 + 1 such values. In total, the number of unsuitable
                                                     p
values x is less than
                                                                                      
                      k                          
                     X        X           N                 X 1             X
                                              + 1   <  k ·  N           +             1 <
                                                                                       
                                          p 2              
                                                                   p 2
                              √
                     i=1 2k<p< M (N +1)                       p>2k
                                                                            √
                                                                          p< M (N +1)
                            X 1            1
                                              
                                                    p                 N     p
                     < kN                −      + k M(N + 1) <           + k M(N + 1).
                            p>2k
                                   p−1 p                               2

If N is big enough then this is less than N, and there exist a suitable choice for x.


Problem 5. Consider a circular necklace with 2013 beads. Each bead can be painted either white or
green. A painting of the necklace is called good, if among any 21 successive beads there is at least one
green bead. Prove that the number of good paintings of the necklace is odd.
    (Two paintings that differ on some beads, but can be obtained from each other by rotating or
flipping the necklace, are counted as different paintings.)
                                            (Proposed by Vsevolod Bykov and Oleksandr Rybak, Kiev)

Solution 1. For k = 0, 1, . . . denote by Nk be the number of good open laces, consisting of k (white
and green) beads in a row, such that among any 21 successive beads there is at least one green bead.
For k ≤ 21 all laces have this property, so Nk = 2k for 0 ≤ k ≤ 20; in particular, N0 is odd, and
N1 , . . . , N20 are even.
    For k ≥ 21, there must be a green bead among the last 21 ones. Suppose that the last green bead
is at the ℓth position; then ℓ ≥ k − 20. The previous ℓ − 1 beads have Nℓ−1 good colorings, and every
such good coloring provides a good lace of length k. Hence,

                             Nk = Nk−1 + Nk−2 + . . . + Nk−21    for k ≥ 21.                          (1)

From (1) we can see that N21 = N0 + . . . + N20 is odd, and N22 = N1 + . . . + N21 is also odd.
   Applying (1) again to the term Nk−1 ,
                                                     
    Nk = Nk−1 + . . . + Nk−21 = Nk−2 + . . . + Nk−22 + Nk−2 + . . . + Nk−21 ≡ Nk−22 (mod 2)

so the sequence of parities in (Nk ) is periodic with period 22. We conclude that

   • Nk is odd if k ≡ 0 (mod 22) or k ≡ 21 (mod 22);

   • Nk is even otherwise.

    Now consider the good circular necklaces of 2013 beads. At a fixed point between two beads cut
each. The resulting open lace may have some consecutive white beads at the two ends, altogether at
most 20. Suppose that there are x white beads at the beginning and y white beads at the end; then
we have x, y ≥ and x + y ≤ 20, and we have a good open lace in the middle, between the first and the
last green beads. That middle lace consist of 2011 − x − y beads. So, for any fixed values of x and y
the number of such cases is N2011−x−y .

                                                    3
                                                              x       2011 − x − y   y


   It is easy to see that from such a good open lace we can reconstruct the original circular lace.
Therefore, the number of good circular necklaces is
 X
      N2011−x−y = N2011 +2N2010 +3N2009 +. . .+21N1991 ≡ N2011 +N2009 +N2007 +. . .+N1991 (mod 2).
x+y≤20


By 91 · 22 − 1 = 2001 the term N2001 is odd, the other terms are all even, so the number of the good
circular necklaces is odd.
Solution 2 (by Yoav Krauz, Israel). There is just one good monochromatic necklace. Let us
count the parity of good necklaces having both colors.
    For each necklace, we define an adjusted necklace, so that at position 0 we have a white bead and
at position 1 we have a green bead. If the necklace is satisfying the condition, it corresponds to itself;
if both beads 0 and 1 are white we rotate it (so that the bead 1 goes to place 0) until bead 1 becomes
green; if bead 1 is green, we rotate it in the opposite direction until the bead 0 will be white. This
procedure is called adjusting, and the place between the white and green bead which are rotated into
places 0 and 1 will be called distinguished place. The interval consisting of the subsequent green beads
after the distinguished place and subsequent white beads before it will be called distinguished interval.
    For each adjusted necklace we have several necklaces corresponding to it, and the number of them
is equal to the length of distinguished interval (the total number of beads in it) minus 1. Since we
count only the parity, we can disregard the adjusted necklaces with even distinguished intervals and
count once each adjusted necklace with odd distinguished interval.
    Now we shall prove that the number of necklaces with odd distinguished intervals is even by grouping
them in pairs. The pairing is the following. If the number of white beads with in the distinguished
interval is even, we turn the last white bead (at the distinguished place) into green. The white interval
remains, since a positive even number minus 1 is still positive. If the number of white beads in the
distinguished interval is odd, we turn the green bead next to the distinguished place into white. The
green interval remains since it was even; the white interval was odd and at most 19 so it will become
even and at most 20, so we still get a good necklace.
    This pairing on good necklaces with distinguished intervals of odd length shows, that the number of
such necklaces is even; hence the total number of all good necklaces using both colors is even. Therefore,
together with monochromatic green necklace, the number of good necklaces is odd.




                                                    4
